% ============================================================================
%  SECCIÓN 5.6: PERSISTENCIA DESPUÉS DEL CIERRE
% ============================================================================

\section{Persistencia despu\'es del cierre}

Uno de los objetivos principales de este avance es garantizar que la
aplicaci\'on conserve los datos cr\'iticos despu\'es del cierre
completo de la aplicaci\'on o del reinicio del dispositivo. A
continuaci\'on se describe c\'omo cada mecanismo de persistencia
opera despu\'es del cierre.

\subsection{Sesi\'on de usuario con AsyncStorage}

Cuando el usuario inicia sesi\'on, Supabase Auth almacena el token
JWT y el refresh token en AsyncStorage. Al reabrir la aplicaci\'on,
el servicio de Supabase recupera autom\'aticamente la sesi\'on
almacenada:

\begin{lstlisting}
// Al iniciar la app: recuperar sesi\'on
supabase.auth.getSession()
  .then(({ data: { session } }) => {
    setSession(session);
    setLoading(false);
  });
\end{lstlisting}

El flujo de persistencia de la sesi\'on es el siguiente:

\begin{enumerate}
  \item El usuario inicia sesi\'on con email y contrase~na.
  \item Supabase retorna un token JWT y un refresh token.
  \item Los tokens se almacenan en AsyncStorage de forma
    autom\'atica.
  \item La aplicaci\'on se cierra o el dispositivo se reinicia.
  \item Al reabrir la app, \textit{getSession()} recupera los
    tokens de AsyncStorage.
  \item Si el token est\'a v\'alido, la sesi\'on se restaura.
  \item Si el token expir\'o, se usa el refresh token para
    obtener uno nuevo.
  \item Si el refresh token tambi\'en expir\'o, el usuario
    debe iniciar sesi\'on nuevamente.
\end{enumerate}

\begin{figure}[H]
\centering
  \figuraAPA{fig:flujo-sesion}{Flujo de persistencia de la sesi\'on}{%
  \begin{tikzpicture}[
    node distance=0.5cm,
    nodo/.style={rectangle, draw=black!60, rounded corners,
      minimum width=4.0cm, minimum height=0.5cm, align=center,
      font=\footnotesize},
    flecha/.style={-Stealth, thick, draw=black!60}
  ]
    \node[nodo] (login) {Inicio de sesi\'on};
    \node[nodo, below=of login] (token) {Token JWT};
    \node[nodo, below=of token] (store) {AsyncStorage};
    \node[nodo, below=of store] (close) {Cierre de app};
    \node[nodo, below=of close] (reopen) {Reabrir app};
    \node[nodo, below=of reopen] (recover) {Recuperar sesi\'on};
    \node[nodo, below=of recover] (valid) {Sesi\'on restaurada};
    \draw[flecha] (login) -- (token);
    \draw[flecha] (token) -- (store);
    \draw[flecha] (store) -- (close);
    \draw[flecha] (close) -- (reopen);
    \draw[flecha] (reopen) -- (recover);
    \draw[flecha] (recover) -- (valid);
  \end{tikzpicture}%
  }{Elaboraci\'on propia.}
\end{figure}

\subsection{Datos booleanos persistentes}

Junto con los tokens, se persisten datos booleanos que reflejan el
estado de una opci\'on o configuraci\'on. El mecanismo de
persistencia local admite este tipo de valor adem\'as de texto,
n\'umeros y listas de cadenas.

En la aplicaci\'on, la sesi\'on del usuario se conserva mediante dos
banderas booleanas provistas por la configuraci\'on de Supabase Auth:
\begin{itemize}
  \item \texttt{autoRefreshToken} -- indica si el token se renueva
    autom\'aticamente antes de expirar.
  \item \texttt{persistSession} -- indica si la sesi\'on debe
    conservarse tras cerrar la aplicaci\'on.
\end{itemize}

Ambas se almacenan de forma persistente y se leen al volver a abrir
la aplicaci\'on, por lo que el estado de la sesi\'on (iniciada o no
iniciada) se conserva correctamente entre cierres.

\subsection{Datos de dominio en Supabase}

Los escenarios, eventos, im\'agenes y favoritos se almacenan
permanentemente en la base de datos de Supabase. Al reabrir la
aplicaci\'on, React Query ejecuta las consultas correspondientes
y obtiene los datos actualizados del servidor:

\begin{lstlisting}
// Al reabrir la app: consultar escenarios
return useQuery({
  queryKey: ['scenarios'],
  queryFn: fetchScenarios,
  staleTime: 1000 * 60 * 5,
});
\end{lstlisting}

Si el usuario ten\'ia favoritos guardados, estos se recuperan
autom\'aticamente desde la tabla \textit{favorites} de Supabase:

\begin{lstlisting}
// Recuperar favoritos del usuario
const { data } = await supabase
  .from('favorites')
  .select(`*, scenarios(*)`)
  .eq('user_id', userId)
  .order('created_at', { ascending: false });
\end{lstlisting}

\subsection{Cach\'e de React Query}

El cach\'e de React Query se pierde al cerrar la aplicaci\'on,
ya que se almacena en memoria. Sin embargo, esto no afecta la
experiencia del usuario porque:

\begin{enumerate}
  \item Las consultas se ejecutan autom\'aticamente al reabrir
    la aplicaci\'on.
  \item Se muestra un indicador de carga mientras se obtienen
    los datos del servidor.
  \item Los datos se cach\'ean nuevamente durante 5 minutos
    para consultas subsiguientes.
\end{enumerate}

Este comportamiento es intencional: React Query est\'a dise~nado
para ser una capa de cach\'e ef\'mera que se regenera a partir del
servidor, no un mecanismo de persistencia.

\subsection{Verificaci\'on de persistencia}

Para verificar que la persistencia funciona correctamente, se
realizaron las siguientes pruebas:

\begin{table}[H]
  \centering
  \caption{Verificaci\'on de persistencia tras cierre}
  \contenidoConFuente{%
    \begin{tablaAPA}
      \begin{tabular}{@{}p{0.7cm}p{3.5cm}p{5.0cm}p{3.0cm}@{}}
        \toprule
        \textbf{ID} & \textbf{Prueba} & \textbf{Resultado} &
        \textbf{Observaci\'on} \\
        \midrule
        V1 & Cerrar y reabrir la app &
        La sesi\'on se mantiene activa sin pedir
        credenciales & AsyncStorage conserva tokens. \\
        \addlinespace
        V2 & Cerrar y reabrir con favoritos &
        Los favoritos aparecen en la lista &
        Supabase conserva datos. \\
        \addlinespace
        V3 & Cerrar y reabrir tras crear escenario &
        El escenario aparece en el cat\'alogo &
        Datos en servidor. \\
        \addlinespace
        V4 & Cerrar y reabrir tras crear evento &
        El evento se muestra en el detalle &
        Datos en servidor. \\
        \addlinespace
        V5 & Desactivar internet y reabrir &
        Se muestran datos del cach\'e anterior &
        React Query cache. \\
        \bottomrule
      \end{tabular}
    \end{tablaAPA}%
  }{Elaboraci\'on propia en base a las pruebas
  realizadas en dispositivo f\'isico.}
\end{table}

\subsection{Pruebas de error: datos vac\'ios o inexistentes}

Adem\'as de verificar que los datos persisten, se comprobaron los
escenarios cr\'iticos en los que la informaci\'on almacenada no existe
o ha sido eliminada. La aplicaci\'on maneja estos casos sin producir
errores y ofrece retroalimentaci\'on al usuario:

\begin{table}[H]
  \centering
  \caption{Pruebas de error sobre persistencia}
  \contenidoConFuente{%
    \begin{tablaAPA}
      \begin{tabular}{@{}p{0.8cm}p{3.6cm}p{2.0cm}p{8.6cm}@{}}
        \toprule
        \textbf{ID} & \textbf{Prueba} & \textbf{Resultado} &
        \textbf{Observaciones} \\
        \midrule
        E1 & Primera ejecuci\'on sin datos &
        Aprobado & Al iniciar la app por primera vez, no hay sesi\'on
        ni favoritos almacenados. La app muestra la pantalla de
        inicio/cat\'alogo y los estados vac\'ios correspondientes, sin
        errores. \\
        \addlinespace
        E2 & Sesiones expiradas o inexistentes &
        Aprobado & Si el refresh token expir\'o o no existe, la app redirige
        al inicio de sesi\'on sin pantallas en blanco ni excepciones. \\
        \addlinespace
        E3 & Favoritos eliminados & Aprobado &
        Tras eliminar todos los favoritos, la lista muestra el estado
        vac\'io (\textit{``A\'un no tienes favoritos''}) en lugar de un error. \\
        \addlinespace
        E4 & Escenarios desactivados & Aprobado &
        Si todos los escenarios est\'an inactivos, el cat\'alogo p\'ublico
        muestra un estado vac\'io con mensaje informativo. \\
        \bottomrule
      \end{tabular}
    \end{tablaAPA}%
  }{Elaboraci\'on propia en base a las pruebas realizadas en dispositivo
  f\'isico.}
\end{table}

\subsection{Comparaci\'on antes y despu\'es del avance}

\begin{table}[H]
  \centering
  \caption{Comparaci\'on de persistencia antes y despu\'es}
  \contenidoConFuente{%
    \begin{tablaAPA}
      \begin{tabular}{@{}p{3.4cm}p{5.0cm}p{6.2cm}@{}}
        \toprule
        \textbf{Aspecto} & \textbf{Antes (PDF 4)} &
        \textbf{Despu\'es (PDF 5)} \\
        \midrule
        Sesi\'on & Se perd\'ia al cerrar la app &
        Se conserva con AsyncStorage. \\
        \addlinespace
        Favoritos & Se perd\'ian al cerrar la app &
        Se conservan en Supabase. \\
        \addlinespace
        Escenarios & Solo lectura (sin CRUD) &
        CRUD completo con persistencia. \\
        \addlinespace
        Eventos & Solo lectura &
        Create y Delete con persistencia. \\
        \addlinespace
        Im\'agenes & Sin gesti\'on &
        Subida y eliminaci\'on con Storage. \\
        \bottomrule
      \end{tabular}
    \end{tablaAPA}%
  }{Elaboraci\'on propia.}
\end{table}
